\chapter{eBPF}

\section{Introduction and overview}

TODO: introduction on eBPF, as the modern solution for dynamic tracing and
observability

TODO: some history of eBPF

TODO: explain how eBPF connects with the rest of the tracing infrastructure

TODO: general overview of the eBPF architecture: system call, verifyier, program
types

TODO: example program

\section{Instruction Set}

eBPF programs are written in the eBPF \textit{instruction set architecture}
(ISA), often referred to simply as BPF ISA. The ISA defines the set of
instructions that an eBPF program uses and that are expected to be supported by
an eBPF runtime. A compiler may be used to generate the BPF instructions from an
higher-level language such as C or Rust to make the development of eBPF program
easier for developers.

For example, the figure below shows C code on the left and an human readable
form of BPF bytecode on the right compiled with clang 23.1.0:

\noindent
\vspace{0.5em}

\begin{minipage}[t]{0.47\textwidth}
\begin{verbatim}
int sum_foo(int a, int b) {
    return foo(a + b);
}
\end{verbatim}
\end{minipage}%
\hfill
\vrule % This creates the vertical dividing bar
\hfill
\begin{minipage}[t]{0.47\textwidth}
\begin{verbatim}
sum_foo(int, int):
        *(u32 *)(r10 - 4) = w1
        *(u32 *)(r10 - 8) = w2
        w1 = *(u32 *)(r10 - 4)
        w2 = *(u32 *)(r10 - 8)
        w1 += w2
        call foo(int)
        exit
\end{verbatim}
\end{minipage}

\noindent
\vspace{0.5em}

Each line on the right corresponds to a BPF assembly instruction. The term
\verb|w1|, \verb|w2| represent the first word of the register number 1 and 2
respectively, and \verb|r10| is the full register number 10.

\subsection{Standardization}

Standardizing the BPF ISA is significantly important to allow compiler
developers, tools developers, eBPF runtime developers and ASIC designers to
agree to a common reference. The formal IETF standardization effort was started
from earlier drafts in October 2023 with RFC 9669 \cite{rfc9669}; the RFC was
finally standardized one year later and updated multiple times since the first
version. All the information presented in this section are referring to the RFC
as last updated in 20 May 2026.

Controversially, the RFC does not standardize the platform-specific application
binary interface (ABI) because a common consensus was not reached, hence it was
preferred to be left out of the standard. This includes various details such as:
which registers are saved across a call, which registers contain the frame
pointer, how large the stack is, and other specifics. Further standardization is
left for "future work"; however, the Linux BPF ABI is well explained in the
kernel documentation and any BPF runtime that wishes to be compatible with the
Linux implementation must follow it.

\subsection{Instructions and encoding}

The standard specifies two instruction encodings: one using 64 bits and another
one using 128 bits. Instructions are assigned to a named conformance group,
which is a set of instructions or other requirements identified with a short
ASCII label, such as "base32". An implementation of the BPF ISA must support the
"base32" conformance group, but may or may not support other groups.

\begin{figure}[htbp]
    \centering
    \begin{bytefield}[bitwidth=1.1em]{32}
        % Header bit numbers (optional)
        \bitheader{0, 7, 15, 31} \\

        % First row
        \bitbox{8}{opcode} & \bitbox{8}{regs} & \bitbox{16}{offset} \\

        % Second row
        \bitbox{32}{imm}
    \end{bytefield}
    \caption{64 bit instruction encoding.}
    \label{fig:instruction-encoding}
\end{figure}

\begin{figure}[htbp]
    \centering
    \begin{bytefield}[bitwidth=1.1em]{32}
        % Optional: Bit header numbers
        \bitheader{0, 7, 15, 31} \\

        % Row 1: opcode, regs, offset
        \bitbox{8}{opcode} & \bitbox{8}{regs} & \bitbox{16}{offset} \\

        % Row 2: imm
        \bitbox{32}{imm} \\

        % Row 3: reserved
        \bitbox{32}{reserved} \\

        % Row 4: next_imm (remember to escape the underscore!)
        \bitbox{32}{next\_imm}
    \end{bytefield}
    \caption{Extended 128-bit instruction encoding.}
    \label{fig:extended-instruction-encoding}
\end{figure}

Figure \ref{fig:instruction-encoding} shows the encoding of 64 bit instructions.
The various field are the following:

\begin{itemize}

  \item \verb|opcode|: the operation to perform, encoded as a \textit{class}
    part and a \textit{specific} part which identifies the instruction on the
    class.

  \item \verb|regs|: source and destination register numbers. Some instructions
    might reuse this fields for other purposes.

  \item \verb|offset|: signed integer offset used with pointer arithmetic, or
    for other purposes.

  \item \verb|imm|: signed integer immediate value.

\end{itemize}

Figure \ref{fig:extended-instruction-encoding} shows the encoding for extended
128 bit instructions. Compared to 64 bit instructions, it adds a \verb|reserved|
32 bit field which is set to 0, and a \verb|next_imm| 32 bit field for a second
signed integer immediate value.

The instruction set is a reduced one in the number of instruction, resembling
RISC-V. It has various instruction classes such as arithmetic operations
\verb|ADD| and \verb|MUL|, jump instructions \verb|JMP| and \verb|JEQ|, load and
store from memory, etc. Additionally, the BTF ISA supports atomic operations for
load and store instructions, which ensure that the operation cannot be
interrupted or corrupted by other accesses to the same memory region by another
BPF program or other actors. Other than atomic \verb|ADD|, \verb|OR|, \verb|AND|
and \verb|XOR|, BPF provides \verb|XCHG| for atomic exchange, and \verb|CMPXCHG|
for atomic compare and exchange. A complete list of the encoding of each
instruction is found in the RFC.

\subsection{Helper Functions}

The RFC introduces the concept of \textit{helper functions} which provide
functionalities implemented by the underlying platform. Notably, helper
functions are not standardized by RFC9669, but they are documented in the Linux
kernel, which will be discussed in a later chapter. To call an helper function,
the jump instruction must be executed with its source register as $0x0$ or
$0x2$, with a platform-specific \verb|ID| identifying a specific function.

\subsection{Maps and platform variables}

RFC9669 defines "maps" as shared memory regions accessible by BPF programs on
some platforms. The particular semantics of a map are not defined by the
standard. Linux supports many kinds of maps which implement various data
structures such as hash maps, arrays, ring buffers etc which provide shared
storage across userspace and kernelspace, and among multiple executions of the
same eBPF program or multiple eBPF programs. Linux maps will be discussed later.

The standard also defines "platform variables" which are memory regions
identified by integer ids exposed by the runtime. They are different than maps
as they are used directly by the program instead of going through helper
functions, conceptually platform variables can be thought of as global
variables.

\subsection{Compiler Support}

Clang 11+ and GCC 12.1+ support compiling C to BPF bytecode. Compilers for other
languages also support the BPF ISA, such as the rust compiler. Atomic operations
which map to atomic instructions can be specified with intrinsics.

\section{The verifier}

% Safe and Unsafe programs

After loading and eBPF program, the BPF code is \textit{just in time} (JIT)
compiled to native machine instructions to be executed as kernel code in a later
moment. This would give eBPF programs the ability to run almost arbitrary code
in kernelspace, with important security implications on confidentiality,
integrity and availability. Without a proper security model, an eBPF program
could potentially read or write arbitrary kernel memory (confidentiality and
integrity issue) or invoke other kernel functions and alter the execution of the
operating system to halt it (availability issue). For these reasons, the eBPF
security model determines strict guardrails and limitations to what a program is
allowed to do. The model discriminates a program as \textit{safe} if it meets
all security requirements, or as \textit{unsafe} otherwise; the user can decide
to disallow the loading of any unsafe programs, which is the default policy on
most Linux distributions. This model is fundamental to the success and wide
adoption of eBPF in production and security-critical environments. Other
security considerations and policies will be discussed in a later section, such
as user permissions and chain of trust.

% Security requirements

The following is a non-exhaustive list of security requirements for safe eBPF
programs:

\begin{itemize}

  \item Programs must always terminate: no infinite loops or infinite recursion
    is allowed (availability issue).

  \item Programs cannot to read or write arbitrary memory (confidentiality,
    integrity issue).

  \item Network programs are not allowed to access memory outside of packet
    bounds (confidentiality, integrity issue).

  \item Programs are not allowed to deadlock (availability issue).

  \item Programs are not allowed to read uninitialized memory (confidentiality
    issue).

\end{itemize}

Notably, the safety requirements are not standardized and may change across
kernel versions. The reason is that the eBPF ecosystem is rapidly evolving,
hence putting effort in standardization is not practical or wanted at the
current state.

% Verifier introduction

The requirements are checked at program load time instead of at runtime, to
trade runtime performance penalties with load-time overhead. If a requirement
cannot be fully checked at load-time, then the eBPF program is rejected or a
method to perform the checks dynamically is provided, although only for a
limited set of functionalities (dynamic pointers). The program that checks the
requirements is known as the eBPF \textit{verifier}: given a BPF program, the
output of the verifier will determine if the program is considered safe or not.
Therefore, the verifier is a critical part of the eBPF architecture and its
implementation has been widely scrutinized. Given the complexity of checking
some of the requirements at load time, the verifier implementation is highly
complex. Just the file \path{$LINUX/kernel/bpf/verifier.c}, which implements
most of the logic, is about $20000$ lines of code long, although solutions to
simplify and modularize it are being discussed by its developers.

The verifier implementation can be divided in two separate moments. First, it
builds a control flow graph of the program, looks for unreachable instructions
and identifies back-edges (loops) among other correctness checks. Second, it
descends every possible execution path in the graph from the first instruction
and looks for unsafe operations, such as out of bounds array accesses and
others. It does so by simulating the execution and keeping track of the state of
each register and the range of possible values it can hold (\textit{value
tracking}). If, during the simulation, a state that was already verified is
encountered during a backward jump, then the verifier considers this path as
safe and does not re-simulate the code again (\textit{pruning}) which greatly
speeds up verification. If any of the safety requirements are not met in any
particular state, then the program is rejected. A more in depth look is now
provided.

\subsection{Initial checks and preparations}

A program's verification starts in \verb|bpf_check()| after the \verb|bpf()|
systemcall is executed with the \verb|BPF_PROG_LOAD| command. Upon verification
entry, the verifier allocates an environment structure of type \verb|struct bpf_verifier_env|.
The environment holds all the state of the verifier for a
particular program, such as jump history, register states (\verb|struct bpf_reg_state|)
and the history of previously explored states, among other information.
Additionally, it contains a log buffer of type \verb|struct bpf_verifier_log|
which logs the state of the program after each instruction;
this is particularly useful to debug verification issues. Since the verification
environment is allocated per-verification, multiple verifications on different
programs can happen simultaneously. Parallel verifications are allowed only if
the program's token (see later) has the \verb|CAP_BPF| capability, since
unprivileged users could exhaust kernel memory and degrade performance by
requesting many verifications. A token is a delegation mechanism where a
privilege process can transfer capabilities to an unprivileged one via a token,
which is then passed to the eBPF program. This will be discussed more in a
following chapter.

After some initializations, the verifier proceeds to perform several checks and
compute various values and structures in the context. Some of these operations
are reported below:

\begin{itemize}

  %\item \verb|bpf_check_btf_info_early()|

  %\item \verb|add_subprog_and_kfunc()|

    \item \verb|check_subprogs()|: check that all jumps are within the same
      subprogram and all subprograms end (a program can contain multiple
      subprograms).

  %\item \verb|bpf_check_btf_info()|

    \item \verb|check_and_resolve_insns()|: check that instructions are valid,
      replace map IDs with actual map pointers

  %\item \verb|bpf_prog_is_offloaded()| which calls
  %   \verb|bpf_prog_offload_verifier_prep()|

    \item \verb|bpf_check_cfg()|:  a non-recursive depth-first-search to detect
      loops in the BPF program. Each instruction, which is represented as a
      \textit{vertex} in a directed graph, can have either a
      \textit{discovered}, \textit{explored}, \textit{fallthrough} or
      \textit{branch} label. The difference between discovered and explored is
      that a discovered instruction did not check all the edges yet, while an
      explored one was fully checked. Each \textit{edge} can either be
      discovered or not. While traversing the graph, the verifier assigns these
      labels to the instructions. If the DAG explores any instruction already
      marked as discovered, a back edge is detected and annotated. This
      information is used during emulation to avoid emulating the entire loop
      with a technique knows as \textit{pruning} which will be explained later.

  %\item \verb|bpf_compute_postorder()|

  %\item \verb|bpf_stack_liveness_init()|

  %\item \verb|check_attach_btf_id()|

  %\item \verb|bpf_compute_const_regs()|

    \item \verb|bpf_prune_dead_branches()|: rewrite conditional branches that
      have constant outcomes into unconditional jumps, so that dead code is not
      checked.

    \item \verb|sort_subprogs_topo()|: sort the subprograms topologically

    \item \verb|bpf_compute_scc()|: : Use a non-recursive adaptation of Tarjan's
      algorithm for strongly connected components computation.


  %\item \verb|bpf_compute_live_registers()|

  %\item \verb|mark_fastcall_patterns()|


\end{itemize}

\subsection{Register types}

To proceed with the next phase of verification, the verifier uses a system of
\textit{register types}: during simulation, each register may have a type
assigned to it. Types provide key information about the type a register holds to
help the verifier make certain decisions. Performing operations on registers may
modify their type, and some operations are not allowed wht certain types. For
example, load and store instructions are allowed only with registers of type
\verb|PTR_TO_CTX|, \verb|PTR_TO_MAP| and \verb|PTR_TO_STACK|.

The set of available types is defined in \verb|enum bpf_reg_type| present in
\path{$LINUX/include/linux/bpf.h}. The following is an extract of the enum
for context:

\begin{verbatim}
enum bpf_reg_type {
  NOT_INIT = 0,     /* nothing was written into register */
  SCALAR_VALUE,     /* reg doesn't contain a valid pointer */
  PTR_TO_CTX,     /* reg points to bpf_context */
  CONST_PTR_TO_MAP,   /* reg points to struct bpf_map */
  PTR_TO_MAP_VALUE,   /* reg points to map element value */
  PTR_TO_MAP_KEY,     /* reg points to a map element key */
  PTR_TO_STACK,     /* reg == frame_pointer + offset */
  PTR_TO_PACKET_META,   /* skb->data - meta_len */
  PTR_TO_PACKET,     /* reg points to skb->data */
  PTR_TO_PACKET_END,   /* skb->data + headlen */
  PTR_TO_FLOW_KEYS,   /* reg points to bpf_flow_keys */
  PTR_TO_SOCKET,     /* reg points to struct bpf_sock */
  PTR_TO_SOCK_COMMON,   /* reg points to sock_common */
  /* ... */
}
\end{verbatim}

\subsection{Program simulation}

The verification proceeds to the next major phase: program simulation. The
simulation descends all possible paths of the program to check for correctness.

The simulation begins with \verb|do_check_main()|. At the very start, the only
valid register to be accessed is \verb|R1| of type \verb|PTR_TO_CTX| which can
be used to read a \verb|struct bpf_context| containing various information about
the program depending on the program type. The verifier keeps the state of the
simulation by allocating a \verb|struct bpf_verifier_state| and \verb|struct bpf_func_state|.
It proceeds to call \verb|do_check| which calls
\verb|do_check_insn| for each instruction of the program; this performs all
checks for an instruction, such as bound checks and type compatibility checks.
While iterating over all the instructions, the verifier keeps a stack of states:
whenever it encounters a conditional jump, it pushes a new state in the stack,
and when an exit condition is reached, the top state is popped and simulated,
until all states are checked (the stack is empty). If all checks were
successful, the verifier calls \verb|do_check_subprogs()| which lazily verifies
all global functions. This is necessary because, during verification, all
functions are assumed to be already verified successfully and checked for
signature match only.

\subsection{Range analsyis}

A powerful feature of the verifier is an abstract interpretation technique known
as \textit{range analysis} \cite{cousot1977abstract}. For each registers, the
verifier keeps track of the range of its possible values (minimum and maximum).
This is important to verify bounded accesses to arrays and other checks. For
example, consider the following C code:

\begin{Verbatim}[numbers=left, numbersep=5pt, frame=single]
unsigned int buff[100];
unsigned int i, x, y; /* initialized somewhere randomly */

if (i < 10) {
  x = buff[i];   /* OK, i is within [0, 10) */
} else {
  y = buff[i];   /* Rejected, i is out of bounds [0, UINT_MAX) */
}
\end{Verbatim}

On line 1, a buffer with 10 elements is created. Every time the buffer is
accessed, the verifier will check that the index used, $i$ in this case, is
always in bounds my tracking the range of values in the register that stores its
value. If the simulation takes the code path where the $if$ statement on line 4
is evaluated to true, then value that $i$ holds must have an upper limit of 10
because of the condition in the $if$ statement. If the verifier cannot ensure
that $i$ does not hold a value smaller than the size of the buffer, then the
program is rejected as shown on line 7. The implementation of range analysis has
been formally verified by \cite{vishwanathan2023verifying}

Additionally, the verifier knows which offsets inside structs can be accessed.
For example, when evaluating the instruction $bpf\_ld\ R0 = *(u32 *)(R6 + 8)$, a
check will be performed to ensure that the address $R6 + 8$ can be loaded, and
if $R6$ is of type \verb|PTR_TO_STACK| then an additional check will be made to
ensure that this address is within stack bounds $[-MAX\_BPF\_STACK, 0)$. It
would also be checked that such stack location was written to previously so that
it would be uninitialized.

\subsection{State pruning}

Some states are saved for the entire verification, for example when performing a
jump to a previous instruction identified by a back-edge during the initial
phase of the verifier. This is used for an optimization called \textit{pruning}:
if the current state of the verification was already visited and verified, then
the verifier prunes the current state, meaning that it considers this execution
path as verified and proceeds to pop the stack to simulate other branches. This
greatly speeds up verification time since some "while" or "for" loops can be
fully pruned instead of emulating every cycle.

After successful verification, the kernel finally proceeds to JIT compile the
program into native assembly. The call stack looks like the following:

\begin{verbatim}
bpf_check
  __bpf_prog_select_runtime
    bpf_prog_jit_compile
      bpf_int_jit_compile (arch/x86/net/bpf_jit_comp.c)
        do_jit
\end{verbatim}

\section{Program Types}

Each eBPF program is of a certain type, which determines what resources and
functions the program has access to, and when the program is executed. The
possible types are enumerated in \verb|enum bpf_prog_type| in
\path{$LINUX/include/uapi/linux/bpf.h}, an extract is reported below:

\begin{verbatim}
enum bpf_prog_type {
  BPF_PROG_TYPE_UNSPEC,
  BPF_PROG_TYPE_SOCKET_FILTER,
  BPF_PROG_TYPE_KPROBE,
  BPF_PROG_TYPE_SCHED_CLS,
  BPF_PROG_TYPE_SCHED_ACT,
  BPF_PROG_TYPE_TRACEPOINT,
  BPF_PROG_TYPE_XDP,
  BPF_PROG_TYPE_PERF_EVENT,
  BPF_PROG_TYPE_CGROUP_SKB,
  /* ... */
};
\end{verbatim}

By convention, loaders based on \textit{libbpf} specify the type of the program
by putting the program's binary code in a specific ELF section of the binary,
which depends upon the program type. For example, a program of type
\verb|BPF_PROG_TYPE_CGROUP_DEVICE| must be put in the \verb|cgroup/dev| ELF
section, and a program of type \verb|BPF_PROG_TYPE_CGROUP_SKB| is put in the
\verb|cgroup/skb| ELF section. This trick allows to specify multiple programs of
different types in the same ELF binary, and it is used by libbpf to provide a
set of convenient APIs for the eBPF program developer.

\subsection{Attach points}

After loading, verifying and JIT compiling the program, it needs to be attached
to some event or other object; the attach point depends on the type of the
program. For example, the program will be attached to a perf event if the
program type is \verb|BPF_PROG_TYPE_PERF_EVENT| or to a function through a
kprobe if the program type is \verb|BPF_PROG_TYPE_KPROBE|. The list of possible
attachment is enumerated in \verb|enum enum bpf_attach_type|. When attaching the
program, an object known as a \textit{link} is created which identifies the
relationship between a program and its attachment point, and manages its
lifetime. The execution of the program depends on its attachment point, for
example a program that is attached to a tracepoint will be executed every time
that tracepoint is fired. The program will also be able to access the
information of the tracepoint if it is attached to it, so that it can implement
specific custom logic to react to the event.

\subsection{Trampolines}

To create certain types of attachments, special trampolines are used, which rely
on the dynamic ftrace infrastructure as explained in a previous chapter, to
allow kernel code to call into BPF programs with practically zero overhead using
code patching. Such special attachments are \verb|BPF_TRACE_FENTRY|,
\verb|BPF_TRACE_FEXIT|, \verb|BPF_MODIFY_RETURN|, \verb|BPF_TRACE_FSESSION|,
\verb|BPF_TRACE_FENTRY_MULTI|, \path{BPF_TRACE_FEXIT_MULTI},
\verb|BPF_TRACE_FSESSIONS_MULTI|. When an eBPF program that uses a trampoline is
linked, the function \verb|bpf_trampoline_update()| is called which creates a
page via \path{arch_prepare_bpf_trampoline()}, setting up the stack environment
for the eBPF program including its arguments, its execution, restoring registers
and finally returning to the target function.

Some example eBPF program types are now analyzed.

\subsection{BPF\_PROG\_TYPE\_KPROBE}

TODO

\subsection{BPF\_PROG\_TYPE\_XDP}

TODO

\subsection{BPF\_PROG\_TYPE\_LSM}

TODO

\subsection{BPF\_PROG\_TYPE\_STRUCT\_OPS}

TODO

\section{Helper functions and Kfuncs}

To perform special operations and to interact with the rest of the operating
system, eBPF are allowed to use special tailored functions for this purpose,
called \textit{helper functions}. Not all programs can use all helper functions:
the eBPF program type determines the set of helpers that can be called; this
ensures that an eBPF program is only allowed to interact with a limited set of
the kernel's subsystems. Helper functions are considered User Space API (UAPI),
meaning that the API is stable and will be kept forward compatible on future
kernel versions. A full list of helpers is documented in the
\textit{bpf-helpers(7)} man page.

Additionally, eBPF programs are allowed to call another set of special
functions: \textit{kfuncs}. These functions can be used to access other
infrastructure in the kernel, but they differ from helper functions as they are
not considered UAPI. This means that an eBPF program that uses kfuncs may not be
compatible with new versions of the kernel, therefore defensive programming
techniques are advised, such as by supporting alternatives in the eBPF program
if a function is not available. Kfuncs are now the preferred method for adding
new functionalities to eBPF programs since they are not expected to be supported
indefinitely: this makes it possible to refactor and improve them with less
effort and reduce code debt. An example of kfunc registration is now analyzed.

The following code defines the kfunc \verb|bpf_list_add()|:

\begin{Verbatim}[numbers=left, numbersep=5pt, frame=single]
__bpf_kfunc_start_defs();

__bpf_kfunc int bpf_list_add(struct bpf_list_head *head, struct bpf_list_node *new,
           struct bpf_list_node *prev__nonown_allowed,
           struct btf_struct_meta *meta, u64 off) {
  /* ... */
}

__bpf_kfunc_end_defs();
\end{Verbatim}

The function needs to be declared between \verb|__bpf_kfunc_start_defs()| and
\path{__bpf_kfunc_end_defs()} on line $1$ and $9$ respectively, which remove
some compiler warnings. The \verb|__bpf_kfunc| specifies some linking options.
They expand with the following code:


\begin{Verbatim}[numbers=left, numbersep=5pt, frame=single]
#define __bpf_kfunc __used __retain __noclone noinline

#define __bpf_kfunc_start_defs()                 \
  __diag_push();                     \
  __diag_ignore_all("-Wmissing-declarations",             \
        "Global kfuncs as their definitions will be in BTF");\
  __diag_ignore_all("-Wmissing-prototypes",             \
        "Global kfuncs as their definitions will be in BTF")

#define __bpf_kfunc_end_defs() __diag_pop()
#define __bpf_hook_start() __bpf_kfunc_start_defs()
#define __bpf_hook_end() __bpf_kfunc_end_defs()
\end{Verbatim}

After declaring the function, it needs to be registered so that the verifier
knows about it, which ensures that only valid kfuncs are called. This is done
with the macro \verb|BTF_ID()| which is used in the following way:

\begin{Verbatim}[numbers=left, numbersep=5pt, frame=single]
BTF_ID(func, bpf_list_add)
\end{Verbatim}

The macro expands to \verb|____BTF_ID|:

\begin{Verbatim}[numbers=left, numbersep=5pt, frame=single]
#define BTF_IDS_SECTION ".BTF_ids"

#define ____BTF_ID(symbol, word)      \
asm(              \
".pushsection " BTF_IDS_SECTION ",\"a\";       \n"  \
".local " #symbol " ;                          \n"  \
".type  " #symbol ", STT_OBJECT;               \n"  \
".size  " #symbol ", 4;                        \n"  \
#symbol ":                                     \n"  \
".zero 4                                       \n"  \
word              \
".popsection;                                  \n");
\end{Verbatim}

For the kfunc, an unique label in the \verb|.BTF_ids| section is created (line
$5$), which is initially filled with 4 bytes of zeros (line $10$) during
compilation and resolved later during the linking phase by a special tool named
\verb|resolve_btfids| which is invoked by the kernel build infrastructure. These
4 byes is filled with the BTF ID of the kfunc, used by the verifier to match it
with the eBPF program. Moreover, additional information about the arguments of
the kfunc can be specified for the verifier via a technique known as
\textit{annotations}, such as whether a null pointer is accepted or that only a
known constant value can be passed.

\section{Maps}

RFC 9669 for the BPF Instruction Set Architecture defines the concept of maps as
"shared memory regions accessible by BPF programs on some platforms". They were
included in the standard because maps are heavily used by the Linux eBPF
implementation as one of the main mechanisms to interact between the eBPF
program and userspace, and among eBPF programs. However, the RFC does not
standardize any specific map interface, leaving this important detail to the
implementation. In this section we examine how maps are implemented in the Linux
kernel. Due to the number of types of maps (more than 30), the specifics of
every maps are not described, their interface is documented in the \verb|bpf|
man page.

Many types of maps exist, specifying different semantics and memory layouts.
These are enumerated in the type \verb|enum bpf_map_type|, and are mostly used
as generic data structures for storage of different types of data. Maps
generally provide an API to implement stacks, queues, ring buffers, arenas etc.
As maps are reference counted, they will stay in memory as long as there are
users of them.

Any map implements a set of operations specified in the \verb|struct bpf_map_ops|
type as shown below from \path{$LINUX/include/linux/bpf.h}:

\begin{Verbatim}[numbers=left, numbersep=5pt, frame=single]
struct bpf_map_ops {
  /* funcs callable from userspace (via syscall) */
  int (*map_alloc_check)(union bpf_attr *attr);
  struct bpf_map *(*map_alloc)(union bpf_attr *attr);
  void (*map_release)(struct bpf_map *map, struct file *map_file);
  void (*map_free)(struct bpf_map *map);
  int (*map_get_next_key)(struct bpf_map *map, void *key, void *next_key);
  void (*map_release_uref)(struct bpf_map *map);
  /* ... */
  /* funcs callable from userspace and from eBPF programs */
  void *(*map_lookup_elem)(struct bpf_map *map, void *key);
  long (*map_update_elem)(struct bpf_map *map, void *key, void *value, u64 flags);
  long (*map_delete_elem)(struct bpf_map *map, void *key);
  long (*map_push_elem)(struct bpf_map *map, void *value, u64 flags);
  long (*map_pop_elem)(struct bpf_map *map, void *value);
  long (*map_peek_elem)(struct bpf_map *map, void *value);
  void *(*map_lookup_percpu_elem)(struct bpf_map *map, void *key, u32 cpu);
  int (*map_get_hash)(struct bpf_map *map);
  /* ... */
}
\end{Verbatim}

The various functions such as \verb|map_alloc()| on line 4 and \verb|map_free()|
on line 6 and many others are called at various moments during the map's
lifetime. The functions \verb|map_lookup_elem()|, \verb|map_update_elem()| and
similar ones are used to modify its contents. The specific implementation of
these functions depends on how the specific map wants to handle the data.

After defining the various operations for a map and defining static storage for
a \verb|struct bpf_maps_ops|, the map is registered with all maps by appending
its name to the \path{$LINUX/include/linux/bpf_types.h} in the macro
\verb|BPF_MAP_TYPE()|. All map entries from this file will be saved in a global
array of maps in the following manner:

\begin{Verbatim}[numbers=left, numbersep=5pt, frame=single]
static const struct bpf_map_ops * const bpf_map_types[] = {
#define BPF_PROG_TYPE(_id, _name, prog_ctx_type, kern_ctx_type)
#define BPF_MAP_TYPE(_id, _ops) \
  [_id] = &_ops,
#define BPF_LINK_TYPE(_id, _name)
#include <linux/bpf_types.h>
#undef BPF_PROG_TYPE
#undef BPF_MAP_TYPE
#undef BPF_LINK_TYPE
};
\end{Verbatim}

The header \verb|linux/bpf_types.h| in line 6 will be included inside the
\verb|bpf_map_types| array, and each \verb|BPF_MAP_TYPE()| macro will be
expanded to place each map's operations in the array based on it's ID on line 4.

To create a map from userspace, the \verb|bpf()| systemcall is used:

\begin{Verbatim}[numbers=left, numbersep=5pt, frame=single]
int bpf_create_map(enum bpf_map_type map_type,
               unsigned int key_size,
               unsigned int value_size,
               unsigned int max_entries)
{
    union bpf_attr attr = {
        .map_type    = map_type,
        .key_size    = key_size,
        .value_size  = value_size,
        .max_entries = max_entries
    };

    return bpf(BPF_MAP_CREATE, &attr, sizeof(attr));
}
\end{Verbatim}

Lookups and other operations are also executed through the \verb|bpf()|
systemcall:

\begin{Verbatim}[numbers=left, numbersep=5pt, frame=single]
int bpf_lookup_elem(int fd, const void *key, void *value)
{
    union bpf_attr attr = {
        .map_fd = fd,
        .key    = ptr_to_u64(key),
        .value  = ptr_to_u64(value),
    };

    return bpf(BPF_MAP_LOOKUP_ELEM, &attr, sizeof(attr));
}
\end{Verbatim}

\section{BTF and CO-RE}

In order to help developers to debug a program, compilers support the generation
of debugging information from the source code. Such information may include:
names, types, function signatures and any language-specific meta data that can
be used to pretty-print or to reference the original source code. On Linux, a
standard format used by userspace programs for encoding debug information is
\textit{DWARF}. For example, both gcc and clang compilers generate debug
information in the DWARF format if the \verb|-g| flag is specified. Many other
compilers use DWARF as the default encoding for debug information, such as the
\textit{go} and \textit{rust} compiler; this is because DWARF was designed to be
a generic and expressive format that can be applied with many languages.
Developers of eBPF programs are also interested in expressive debug information,
however, the DWARF format was considered too complex and verbose due to its
design decisions; inspired by the minimalistic and compact design of
\textit{Compact C Type Format} (CTF) developed by Sun Microsystems in Solaris 9,
a custom format was developed specifically for eBPF: the \textit{BPF Type
Format} (BTF). BTF encodes information for all the necessary types, such as
unions, structs, enums, typesefs, etc, and can be used to pretty-print eBPF
maps, functions and structs, among other things. Compilers usually embeds them
in the \verb|.BTF| ELF section of the eBPF program. The details of the format
are documented in the Linux kernel documentation.

For example, consider the following code:

\begin{Verbatim}[numbers=left, numbersep=5pt, frame=single]
struct foo {
  int a;
  int b;
  unsigned c:15;
} __attribute__((preserve_access_index));

enum bar { U, V };
\end{Verbatim}

Given the code above, clang will generate the following BTF information, qpshown
in an human readable form:

\begin{Verbatim}[numbers=left, numbersep=5pt, frame=single]
 [2] STRUCT 'foo' size=8 vlen=2
        'a' type_id=3 bits_offset=0
        'b' type_id=3 bits_offset=32
        'c' type_id=4 bits_offset=64 bitfield_size=15
 [3] INT 'int' size=4 bits_offset=0 nr_bits=32 encoding=SIGNED
 [4] INT 'unsigned int' size=4 bits_offset=0 nr_bits=32 encoding=(none)
 ...
 [16] ENUM 'bar' encoding=UNSIGNED size=4 vlen=2
        'U' val=0
        'V' val=1
\end{Verbatim}

Form the clang output we can see that the name and the size of the struct is
saved in line 1, along with the information of its members and their offsets.
Each type is assigned an unique number \verb|type_id| which is used to reference
the type.

Indeed each type's debug information is encoded in the following compact format,
from \path{$LINUX/include/uapi/linux/btf.h}:

\begin{Verbatim}[numbers=left, numbersep=5pt, frame=single]
struct btf_type {
	__u32 name_off;
	/* "info" bits arrangement
	 * bits  0-23: vlen (e.g. # of struct's members)
	 * bits 24-30: kind (e.g. int, ptr, array...etc)
	 * bit     31: kind_flag, currently used by
	 *             struct, union, enum, fwd, enum64,
	 *             decl_tag and type_tag
	 */
	__u32 info;
	/* "size" is used by INT, ENUM, STRUCT, UNION, DATASEC and ENUM64.
	 * "size" tells the size of the type it is describing.
	 *
	 * "type" is used by PTR, TYPEDEF, VOLATILE, CONST, RESTRICT,
	 * FUNC, FUNC_PROTO, VAR, DECL_TAG and TYPE_TAG.
	 * "type" is a type_id referring to another type.
	 */
	union {
		__u32 size;
		__u32 type;
	};
};
\end{Verbatim}

BTF provides useful metadata to debug eBPF programs, but another important
use-case was found for the same information, which is not discussed.

While the eBPF sybsystem provides a stable sandbox for eBPF programs, an issue
appears when we have programs that interact with the rest of the kernel's types
and structures, such as programs that attach to kprobes and traceporints. The
kernel's internal API is famously unstable, which allows for fast development,
code refactoring and decreases technical debth over time. However, this decision
makes it difficult to distribute external code that is loaded in the kernel;
kernel modules have a strict requirement that they need to be compiled for the
same kernel version as the kernel they are loaded with for this exact reason,
which is unconvenient. Due to the popularity and usefullness of eBPF a new
system was designed.

A first solution to this problem was attempted by the \textit{BPF Compiler
Collection} (BCC). The approac was to distribute C source code directly, then
compiled on the target machine. This technically worked but it required the
entire compilation toolchain (clang/LLVM) to be installed on the host machine.
Since LLVM is considered an heavy dependency, a new solution had to be designed.

The new approach was called \textit{Compile Once, Run Everywhere} (CO-RE) and is
based on relocations, like the ones used when linking a program. The key idea is
to mark memory accesses as relocatable items, which are resolved by the loader
program using BTF information before loading the program. This can be performed
by a self contained executable without requiring heavy dependencies. In order to
resolve the relocations, the loader needs to know the BPF bytecode of the
program, its relocation records, and the BTF information of all kernel types and
symbols. This last data is exposed to userspace via the files
\path{$SYSFS/kernel/btf/vmlinux} and \path{$SYSFS/kernel/btf/$MODULE} for
dynamically loaded modules. With this approach, if the address of a member in a
struct changes, then the relocation will take care of setting the right offset
before the program is loaded, leveraging BTF information in order to make the
eBPF program portable across kernel versions. The main steps of this operation
are now described. Before loading a program, the loader reads the program's
relocation records which the compiler generated in the \verb|.BTF.ext| section.
Then, it checks the current kernel's BTF information from the files in sysfs,
and resolves the relocations using the BTF information. It finally loads the
final progam via the bpf systemcall. The kernel is completely CO-RE agnostic
since the relocations are done in userspace by libbpf.

\section{Security Considerations}

TODO

Status: done researching this section, needs to be written

\section{Use Cases and modern ecosystem}

TODO

Status: done researching this section, needs to be written
